% Generated draft: requires mathematical and source review.
\documentclass[11pt]{article}
\usepackage[margin=25mm]{geometry}
\usepackage{fontspec}
\setmainfont{Latin Modern Roman}
\usepackage{amsmath,amssymb,mathtools,unicode-math}
\setmathfont{Latin Modern Math}
\usepackage{enumitem,needspace,xurl,hyperref}
\setlength{\parindent}{0pt}
\setlength{\parskip}{5pt}
\setlength{\emergencystretch}{3em}
\newcommand{\N}{\mathbb{N}}
\newcommand{\Z}{\mathbb{Z}}
\newcommand{\Q}{\mathbb{Q}}
\newcommand{\R}{\mathbb{R}}
\newcommand{\C}{\mathbb{C}}
\newcommand{\F}{\mathbb{F}}
\newcommand{\A}{\mathbb{A}}
\newcommand{\PP}{\mathbb{P}}
\newcommand{\E}{\mathbb{E}}
\newcommand{\eps}{\varepsilon}
\newcommand{\dd}{\,\mathrm{d}}
\begin{document}
\section*{UnsolvedMath ID: 30001296}
\textbf{Source problem code:} catalogue entry
\subsection{Definitions and mathematical statement}
Let \(f_0\in\Q[x_0,\ldots,x_{n+1}]\) be homogeneous of degree \(d\), and put
\[
f_1=x_0^d+\cdots+x_{n+1}^d.
\]
Consider the pencil of projective hypersurfaces of degree \(d\) and dimension \(n\),
\[
X_t:\quad (1-t)f_0+tf_1=0.
\]
Assume that this pencil is generically smooth; the member at \(t=0\) is allowed to be singular.

For a prime \(p\) satisfying the source's unspecified suitability conditions, let \(V\) denote the associated finite-dimensional limiting cohomology space over \(\Q_p\). Let \(F\) and \(N\) be its Frobenius and nilpotent monodromy operators. The supplied record specifies matrices with
\[
F\in\operatorname{Mat}_{r}(\Q_p),\qquad N\in\operatorname{Mat}_{r}(\Q),\qquad NF=pFN,
\]
where \(r=\dim_{\Q_p}V\), and the rational matrix \(N\) acts on \(V\) through \(\Q\subset\Q_p\). Here nilpotence means that \(N^m=0\) for some positive integer \(m\). Define
\[
K=\ker(N:V\longrightarrow V).
\]
The relation \(NF=pFN\) makes \(K\) invariant under \(F\), so the polynomial
\[
P(T)=\det\bigl(I_K-T(F|_K)\bigr)\in\Q_p[T]
\]
is defined.

The supplied conjecture asks whether, for every such pencil and every prime satisfying the intended suitability conditions,
\[
P(T)\in\Z[T],
\]
and whether every complex reciprocal root \(\alpha\) of this integer polynomial satisfies
\[
|\alpha|=p^{i/2}\qquad\text{for some }i\in\Z\text{ with }0\leq i\leq n.
\]
Here a complex reciprocal root means a nonzero \(\alpha\in\C\) such that \(P(\alpha^{-1})=0\), after regarding the asserted integer polynomial as a complex polynomial. The integer \(i\) may depend on the reciprocal root.

The formulation remains conditional: the supplied record does not define the limiting cohomology construction, its cohomological degree, or the suitability conditions on \(p\). It also does not specify whether the cohomology is full or primitive, whether a base change is required, or which Frobenius normalization is intended. These choices must be recovered from the cited report before the conjecture has a fully specified scope. The projective interpretation of the homogeneous hypersurfaces is made explicit here.
\subsection{Short English statement}
For a generically smooth pencil joining a rational homogeneous hypersurface to the Fermat hypersurface, the conjecture predicts that Frobenius on the kernel of limiting monodromy has an integer characteristic polynomial in reciprocal form, whose complex reciprocal roots have absolute values \(p^{i/2}\) with \(0\leq i\leq n\). The precise limiting cohomology and admissible primes remain unspecified in the supplied record.
\subsection{Sources}
\begin{itemize}
\item Supplied local UnsolvedMath registry record, permanent ID \(30001296\), source problem code \texttt{catalogue entry}, title: \emph{Frobenius Polynomials of Degeneration Pencils}; collection: \emph{Oberwolfach Reports Open Problems}. No individual proposer is identified; the supplied proposed year is 2009.
\item Supplied report citation: \emph{Explicit Methods in Number Theory (2010)}, Oberwolfach Reports, DOI 10.4171/owr/2009/33, pp.\ 1843--1920. Source URL and supplied literature-source link labelled ``OWR report'': \url{https://doi.org/10.4171/owr/2009/33}. The record describes these reports as accounts of mathematical workshops and their open-problem discussions.
\item Supplied literature assessment, dated exactly \texttt{2026-08-21}: status ``open''; the report is said to state the integrality and weight conjecture and give examples, while the supplied review reports finding no general proof. This is supplied dated curation, not an independently verified literature assessment. The listed source checked is the OWR report linked above.
\item The supplied statement status is \texttt{corrected\_verified}, with a supplied assertion that the statement was checked against the official report and preserves its hypotheses, notation, and open target. That verification assertion is recorded as provenance and has not been independently confirmed here.
\end{itemize}
\end{document}
